\documentclass{beamer}
\usetheme{Madrid}
\usecolortheme{default}

\usepackage{amsmath,amssymb}
\usepackage{algorithm,algorithmicx,algpseudocode}

\title{DP-GCD:\\High-Dimensional Private Empirical Risk Minimization\\by Greedy Coordinate Descent}
\author{Pierre Cornilleau, Amine Belkasmi}
\institute{Differential Privacy in ML project}
\date{March 2026}

\begin{document}

\frame{\titlepage}

\begin{frame}{Outline}
\tableofcontents
\end{frame}

\section{Motivation}

\begin{frame}{The Privacy-Utility Dilemma}
\begin{columns}[T]
\begin{column}{0.5\textwidth}
\textbf{Why Privacy Matters}
\begin{itemize}
\item Medical records
\item Financial data
\item Personal information
\item Legal requirements (GDPR)
\end{itemize}

\vspace{1em}

\textbf{The Problem}
\begin{itemize}
\item Models leak training data
\item Membership inference attacks
\item Model inversion attacks
\end{itemize}
\end{column}

\begin{column}{0.5\textwidth}
% Placeholder for privacy-utility trade-off figure
\begin{center}
\textit{[Privacy-utility trade-off illustration]}
\end{center}
\end{column}

\vspace{0.5em}

\textbf{Differential Privacy (DP)} provides rigorous guarantees but at a cost...
\end{columns}

\vspace{0.5em}

\alert{Course connection:} Fundamental privacy-utility trade-off from lecture
\end{frame}

\begin{frame}{The Curse of Dimensionality in DP-ERM}
\textbf{Standard Problem:}
\[
\min_{w \in \mathbb{R}^p} f(w) = \frac{1}{n} \sum_{i=1}^{n} \ell(w, d_i)
\quad \text{with DP}
\]

\vspace{1em}

\textbf{Known Results} (from course + lower bounds lecture):
\begin{itemize}
\item Bassily et al.\ (2014): DP-ERM utility degrades polynomially with dimension $p$
\item For high-dimensional problems $p \gg n$: catastrophic utility loss
\item Lower bound: $\Omega(\sqrt{p/n})$ for convex, $\ell_2$-bounded problems
\end{itemize}

\vspace{1em}

\begin{block}{The Challenge}
Can we achieve \textbf{logarithmic} dependence on $p$ instead of polynomial?
\end{block}
\end{frame}

\begin{frame}{Prior Approaches and Limitations}
\begin{table}
\begin{tabular}{l|l|l}
\textbf{Approach} & \textbf{Dimension Dependence} & \textbf{Limitation} \\ \hline
DP-SGD & $O(p / n^{2/3})$ & Polynomial in $p$ \\
DP-Frank-Wolfe & $O(\sqrt{p} / n^{2/3})$ & Requires $\ell_1$ constraint \\
Sparse regression & $O(\log p)$ & Needs known sparsity \\
Public data methods & $O(\log p)$ & Requires public data \\
\end{tabular}
\end{table}

\vspace{1em}

\begin{alertblock}{Gap}
No algorithm achieves $\log p$ dependence for unconstrained problems without prior structure knowledge
\end{alertblock}

\vspace{0.5em}

\alert{Course connection:} Structure exploitation to beat worst-case bounds
\end{frame}

\section{Background}

\begin{frame}{Report-Noisy-Max: The Key Mechanism}
\textbf{Exponential Mechanism} (from the course):
\begin{itemize}
\item Utility function $u(x, y)$, sensitivity $\Delta_u$
\item Sample from $\Pr[M(x) = y] \propto \exp\left( \frac{u(x,y)}{2\Delta_u} \right)$
\end{itemize}

\vspace{0.8em}

\textbf{Report-Noisy-Max} as Special Case:
\begin{itemize}
\item Given scores $q_1, \ldots, q_p$: $\tilde{j} = \arg\max_j \{ q_j + \mathrm{Lap}(\lambda_j) \}$
\item Reveals only \textbf{1 index} instead of $p$ scalar values
\item \alert{Exponential privacy savings} compared to privatizing all $p$ values!
\end{itemize}

\vspace{0.8em}

\textbf{Why this matters for DP-GCD:}
\begin{itemize}
\item DP-SGD: privatize all $p$ gradient coordinates $\Rightarrow$ noise $\propto \sqrt{p}$
\item DP-GCD: Report-Noisy-Max for coordinate selection + privatize only 1 coordinate
\item Noise scales as $\sqrt{T}$ not $\sqrt{p}$ when $T \ll p$ $\Rightarrow$ massive savings!
\end{itemize}
\end{frame}

\begin{frame}{Coordinate Descent Basics}
\textbf{Standard Gradient Descent:}
\[
w_{t+1} = w_t - \nabla f(w_t) \quad \text{(update all $p$ coordinates)}
\]

\vspace{1em}

\textbf{Coordinate Descent:}
\[
w_{t+1} = w_t - \nabla_j f(w_t) \cdot e_j \quad \text{(update only coordinate $j$)}
\]

\vspace{1em}

\begin{columns}[T]
\begin{column}{0.5\textwidth}
\textbf{Random Selection}
\begin{itemize}
\item Pick $j$ uniformly at random
\item Simple, unbiased
\item May waste updates
\end{itemize}
\end{column}

\begin{column}{0.5\textwidth}
\textbf{Greedy Selection}
\begin{itemize}
\item Pick $j$ with largest $|\nabla_j f|$
\item Faster convergence
\item Exploits structure
\end{itemize}
\end{column}
\end{columns}

\vspace{0.5em}

\alert{Course connection:} Optimization algorithms lecture—first-order methods
\end{frame}

\section{The DP-GCD Algorithm}

\begin{frame}{Algorithm Overview}
\begin{algorithm}[H]
\caption{DP-GCD (Simplified)}
\begin{algorithmic}[1]
\State \textbf{Input:} $w_0 \in \mathbb{R}^p$, iterations $T$, noise scales $\lambda_j, \lambda'_j$
\For{$t = 0$ to $T-1$}
    \State Compute gradient $g = \nabla f(w_t)$
    \State \textbf{Greedy selection (private):} $j_t \gets \arg\max_{j \in [p]} \left\{ |g_j|/\sqrt{M_j} + \mathrm{Lap}(\lambda_j) \right\}$
    \State \textbf{Noisy gradient coordinate:} $\tilde{g}_{j_t} \gets g_{j_t} + \mathrm{Lap}(\lambda'_{j_t})$
    \State \textbf{Single coordinate update:} $w_{t+1} \gets w_t - \gamma_{j_t} \tilde{g}_{j_t} e_{j_t}$
\EndFor
\State \textbf{Return} $w_T$
\end{algorithmic}
\end{algorithm}

\vspace{0.5em}

\textbf{Key idea:} Focus privacy budget on most informative coordinates

\vspace{0.3em}

\alert{Course connection:} Combining mechanism design with optimization
\end{frame}

\begin{frame}{Why Greedy Selection Helps Privacy}
\textbf{DP-SGD approach:}
\begin{itemize}
\item Privatize all $p$ gradient coordinates $\Rightarrow$ noise scales as $\sqrt{p}$
\item Wasted noise on irrelevant coordinates
\end{itemize}

\vspace{1em}

\textbf{DP-GCD approach:}
\begin{itemize}
\item Report-Noisy-Max: Select coordinate privately (reveals only index)
\item Privatize only 1 coordinate per iteration
\item Noise scales as $\sqrt{T}$, not $\sqrt{p}$
\item When $T \ll p$: massive savings!
\end{itemize}

\vspace{1em}

\begin{block}{Intuition}
For sparse/quasi-sparse problems, only $s \ll p$ coordinates matter. DP-GCD automatically finds and updates them.
\end{block}

\vspace{0.3em}

\alert{Course connection:} Sensitivity analysis—coordinate-wise sensitivity $\Delta_j = 2L_j/n$ enables fine-grained noise
\end{frame}

\begin{frame}{Privacy Guarantees}
\begin{theorem}[Privacy, Theorem 4.2]
With noise scales
\[
\lambda_j = \lambda'_j = \frac{8 L_j}{n} \frac{\sqrt{T \log(1/\delta)}}{\varepsilon},
\]
DP-GCD is $(\varepsilon, \delta)$-differentially private.
\end{theorem}

\vspace{0.5em}

\textbf{Proof sketch} (using course concepts):
\begin{enumerate}
\item Each iteration: 2 data accesses (selection + gradient)
\item \textbf{Sensitivity:} Coordinate $j$ gradient has $\ell_1$-sensitivity $\Delta_j = 2L_j/n$
\item \textbf{Laplace mechanism:} Adding $\mathrm{Lap}(\Delta_j/\varepsilon_t)$ gives $\varepsilon_t$-DP
\item \textbf{Advanced composition:} Over $2T$ mechanisms with $\varepsilon_t = \varepsilon/(2\sqrt{2T\log(1/\delta)})$ yields $(\varepsilon, \delta)$-DP
\end{enumerate}

\vspace{0.5em}

Solving for $\lambda_j$:
\[
\lambda_j = \frac{\Delta_j}{\varepsilon_t} = \frac{2L_j/n}{\varepsilon/(2\sqrt{2T\log(1/\delta)})} = \frac{8L_j\sqrt{T\log(1/\delta)}}{n\varepsilon}
\]

\vspace{0.3em}

\textbf{Key insight:} Advanced composition theorem gives $\sqrt{T}$ dependence (basic composition would give linear $T$)
\end{frame}

\section{Theoretical Results}

\begin{frame}{Main Utility Results}
\begin{theorem}[Utility, Theorem 4.4, Simplified]
Let $R_{M,1} = \|w_0 - w^*\|_{M,1}$. With high probability:

\vspace{0.5em}

\textbf{Convex case:}
\[
f(w_{\text{priv}}) - f^* = \widetilde{O}\!\left( R_{M,1}^{4/3} \, L^{2/3} \, \varepsilon^{-1/3} \left(\frac{\log p}{n}\right)^{2/3} \right)
\]

\textbf{Strongly convex case:}
\[
f(w_{\text{priv}}) - f^* = \widetilde{O}\!\left( \frac{L^2 \, \varepsilon^{-2} \log p}{\mu_{M,1} \, n^2} \right)
\]
\end{theorem}

\vspace{1em}

\begin{alertblock}{Key Achievement}
\textbf{Logarithmic} dependence on $p$ (not polynomial)!
\end{alertblock}

\vspace{0.3em}

\alert{Course connection:} Sample complexity—$n^{-2}$ rate matches non-private optimal (up to $\log$ factors)
\end{frame}

\begin{frame}{Proof Idea: Strongly Convex Case}
\textbf{Step 1: Noisy descent lemma}
\begin{itemize}
\item By smoothness: $f(w_{t+1}) \leq f(w_t) - \frac{1}{2M_j}|\nabla_j f(w_t)|^2 + O(\text{noise}^2)$
\end{itemize}

\vspace{0.5em}

\textbf{Step 2: Greedy selection advantage}
\begin{itemize}
\item Greedy rule ensures: $\frac{1}{M_j}|\nabla_j f(w_t)|^2 \geq \frac{\|\nabla f(w_t)\|_2^2}{4\sum_j M_j} - O(\text{noise})$
\end{itemize}

\vspace{0.5em}

\textbf{Step 3: Strong convexity}
\begin{itemize}
\item $\|\nabla f(w_t)\|^2 \geq 2\mu_{M,1}(f(w_t) - f^*)$
\end{itemize}

\vspace{0.5em}

\textbf{Step 4: Recursion}
\begin{itemize}
\item $f(w_{t+1}) - f^* \leq (1 - \Theta(\mu))(f(w_t) - f^*) + O(\lambda^2)$ where $\lambda \propto \varepsilon^{-1}\sqrt{T}/n$
\item After $T \sim \mu^{-1}$ iterations: exponential convergence + noise accumulation
\item Balancing gives $O(L^2 \varepsilon^{-2} \log p / n^2)$
\end{itemize}

\vspace{0.3em}

\alert{Course connection:} Contraction analysis from optimization lecture
\end{frame}

\begin{frame}{Comparison with Prior Work}
\begin{table}
\begin{tabular}{l|c|c|c}
\textbf{Algorithm} & \textbf{Convex} & \textbf{Strongly Convex} & \textbf{Constraint} \\ \hline
DP-SGD & $O(p/n^{2/3})$ & $O(p/n^2)$ & None \\
DP-FW & $O(\sqrt{p}/n^{2/3})$ & $O(p^{4/3}/n^{4/3})$ & $\ell_1$ ball \\
\rowcolor{yellow!30}
DP-GCD & $\widetilde{O}(\log p/n^{2/3})$ & $\widetilde{O}(\log p/n^2)$ & None \\
\end{tabular}
\end{table}

\vspace{1em}

\textbf{Breakthrough:}
\begin{itemize}
\item Matches DP-FW without $\ell_1$ constraint
\item First to achieve $\widetilde{O}(\log p / n^2)$ for strongly convex
\item Works on general unconstrained problems
\end{itemize}
\end{frame}

\begin{frame}{Adaptation to Quasi-Sparsity}
\begin{definition}[$s$-Quasi-Sparsity]
Vector $w$ is $(s, \tau)$-quasi-sparse if at most $s$ entries exceed $\tau$ in magnitude.
\end{definition}

\vspace{1em}

\begin{theorem}[Theorem 4.7, Simplified]
For strongly convex problems with quasi-sparse solution, if iterates remain $s$-sparse:
\[
f(w_T) - f^* = \widetilde{O}\!\left( \frac{L^2 \, s^2 \, \varepsilon^{-2} \log p}{\mu \, n^2} \right)
\]
\end{theorem}

\vspace{1em}

\textbf{Impact:} Dimension dependence reduced from $\log p$ to $s^2 \log p$

\vspace{0.5em}

\alert{Crucial:} DP-GCD adapts automatically without knowing $s$ a priori!

\vspace{0.3em}

\alert{Course connection:} Effective dimension vs.\ ambient dimension (high-dim stats lecture)
\end{frame}

\begin{frame}{How DP-GCD Exploits Structure}
\begin{center}
\textit{[Illustration: quasi-sparse weight vector with few large coordinates]}
\end{center}

\textbf{DP-GCD behavior:}
\begin{itemize}
\item Greedy selection focuses on large coordinates early
\item Sparse iterates when initialized at $w_0 = 0$ (``screening effect'')
\item Converges in $T \sim s$ iterations (not $T \sim p$)
\end{itemize}
\end{frame}

\section{Experiments}

\begin{frame}{Experimental Setup}
\begin{table}
\begin{tabular}{l|c|c|l}
\textbf{Dataset} & $n$ & $p$ & \textbf{Notes} \\ \hline
log1, log2 & 1,000 & 100 & Synthetic, quasi-sparse \\
square & 1,000 & 1,000 & Synthetic, 10 non-zeros \\
california & 20,640 & 8 & Real, low-dimensional \\
madelon & 2,600 & 501 & Real \\
dorothea & 800 & 5,000 & Synthetic (high-dim) \\
\end{tabular}
\end{table}

\vspace{1em}

\textbf{Algorithms compared:}
\begin{itemize}
\item DP-SGD: Batch size 1, Gaussian noise
\item DP-CD: Random coordinate selection
\item DP-GCD: Greedy coordinate selection (proposed)
\end{itemize}

\vspace{0.5em}

\textbf{Privacy budget:} $\varepsilon = 1$, $\delta = 1/n^2$ (all experiments)

\vspace{0.5em}

\textbf{Metric:} Relative error $(f(w_t) - f^*)/f^*$ vs.\ passes on data
\end{frame}

\begin{frame}{Key Experimental Results}
\begin{columns}[T]
\begin{column}{0.5\textwidth}
\textbf{High-dim sparse (square, dorothea):}
\begin{itemize}
\item DP-GCD: \alert{Only algorithm that converges}
\item DP-SGD/DP-CD: Cannot decrease objective ($p \gg n$)
\item Dramatic advantage
\end{itemize}

\vspace{1em}

\textbf{Quasi-sparse (log2, madelon):}
\begin{itemize}
\item DP-GCD: 2--3$\times$ faster convergence
\item Reaches $10^{-3}$ error vs.\ $10^{-2}$ for baselines
\end{itemize}
\end{column}

\begin{column}{0.5\textwidth}
\textbf{Low-dimensional (california):}
\begin{itemize}
\item DP-GCD: Comparable to DP-SGD
\item Greedy overhead not justified ($p = 8$)
\item No sparsity to exploit
\end{itemize}

\vspace{1em}

\textbf{Computational cost:}
\begin{itemize}
\item 1 DP-GCD iter = full gradient
\item Equivalent to $p$ DP-CD iters
\item Justified by better utility per pass
\end{itemize}
\end{column}
\end{columns}

\vspace{1em}

\begin{alertblock}{Takeaway}
DP-GCD shines on high-dimensional problems with quasi-sparse structure
\end{alertblock}
\end{frame}

\begin{frame}{Visualization: log2 Results}
\begin{center}
\textit{[Insert your log2 convergence plot here]}
\end{center}

\vspace{-0.5em}

\textbf{Observation:} DP-GCD (blue) converges much faster than DP-CD (red) and DP-SGD (green) on quasi-sparse logistic regression problem.
\end{frame}

\section{Critical Analysis}

\begin{frame}{Strengths}
\begin{enumerate}
\item \textbf{No domain constraints}
\begin{itemize}
\item Works on unconstrained problems (unlike DP-Frank-Wolfe)
\end{itemize}

\item \textbf{Automatic adaptation}
\begin{itemize}
\item Exploits quasi-sparsity without prior knowledge of $s$ or $\tau$
\end{itemize}

\item \textbf{Logarithmic dimension dependence}
\begin{itemize}
\item Achieves $\widetilde{O}(\log p)$ in favorable settings
\end{itemize}

\item \textbf{Coordinate-wise adaptivity}
\begin{itemize}
\item Uses fine-grained Lipschitz constants $L_j$ for noise calibration
\end{itemize}

\item \textbf{First optimal strongly-convex rates}
\begin{itemize}
\item $\widetilde{O}(\log p / n^2)$ without reduction tricks
\end{itemize}

\item \textbf{Strong empirical validation}
\begin{itemize}
\item Dramatic improvements on dorothea ($p \sim 5{,}000$)
\end{itemize}
\end{enumerate}
\end{frame}

\begin{frame}{Limitations}
\begin{enumerate}
\item \textbf{Computational cost}
\begin{itemize}
\item Full gradient per iteration: $O(np)$ cost
\item Wall-clock time can be higher than randomized methods
\item Only justified when utility per pass matters more than time
\end{itemize}

\item \textbf{Proximal extension lacks theory}
\begin{itemize}
\item Proposes proximal variant for $\ell_1$ regularization
\item No utility analysis provided
\item Greedy proximal CD is challenging even without privacy
\end{itemize}

\item \textbf{Assumption requirements}
\begin{itemize}
\item Needs coordinate-wise smoothness/Lipschitzness constants
\item Lower bound on initial gap $f(w_0) - f^*$
\end{itemize}

\item \textbf{Not optimal in all regimes}
\begin{itemize}
\item When $\ell_2$-norm quantities bounded: doesn't match DP-SGD
\item Interpolates but doesn't dominate everywhere
\end{itemize}
\end{enumerate}
\end{frame}

\begin{frame}{Limitations (continued)}
\begin{enumerate}
\setcounter{enumi}{4}
\item \textbf{Privacy budget overhead}
\begin{itemize}
\item Report-Noisy-Max adds noise for coordinate selection
\item In low-dim: overhead outweighs benefits
\end{itemize}

\item \textbf{Limited to smooth objectives}
\begin{itemize}
\item Main theory assumes smooth losses
\item Extensions to non-smooth/constrained settings remain open
\end{itemize}

\item \textbf{Gap between theory and practice}
\begin{itemize}
\item Theoretical noise scales may be too conservative
\item Experiments use tuned hyperparameters (not theoretical formulas)
\end{itemize}
\end{enumerate}

\vspace{1em}

\begin{block}{When to Use DP-GCD}
High-dimensional problems with suspected quasi-sparsity, where privacy-utility trade-off matters more than wall-clock time.
\end{block}
\end{frame}

\section{Conclusion}

\begin{frame}{Summary}
\textbf{Main Contributions:}
\begin{enumerate}
\item \textbf{DP-GCD algorithm:} Greedy coordinate descent with Report-Noisy-Max
\item \textbf{Theoretical guarantees:} $\widetilde{O}(\log p)$ utility for unconstrained problems
\item \textbf{Automatic adaptation:} Exploits quasi-sparsity without prior knowledge
\item \textbf{Empirical validation:} Dramatic improvements on high-dim sparse problems
\end{enumerate}

\vspace{1em}

\begin{block}{Key Insight}
By focusing privacy budget on most informative coordinates, DP-GCD achieves logarithmic dimension dependence and automatically adapts to problem structure.
\end{block}

\vspace{1em}

\textbf{Impact:} Opens door to practical private learning in high dimensions (genomics, NLP, high-dim regression).
\end{frame}

\begin{frame}{Open Questions}
\begin{enumerate}
\item \textbf{Proximal analysis:} Can we rigorously analyze the proximal variant for non-smooth regularizers?

\item \textbf{Optimal algorithm:} Does there exist an algorithm achieving optimal utility in all regimes (both $\ell_1$ and $\ell_2$ bounded)?

\item \textbf{Computational efficiency:} Can we reduce per-iteration cost while preserving privacy-utility benefits?

\item \textbf{Extensions:} How does DP-GCD perform on non-convex problems (neural networks)?

\item \textbf{Fairness implications:} Does sparsity exploitation introduce fairness concerns (favoring represented groups)?
\end{enumerate}

\vspace{1em}

\begin{alertblock}{Research Direction}
Toward adaptive DP algorithms that automatically exploit problem structure without sacrificing worst-case guarantees.
\end{alertblock}
\end{frame}

\begin{frame}[standout]
Thank You!

\vspace{2em}

Questions?

\vspace{2em}

\small
Mangold, P., Bellet, A., Salmon, J., and Tommasi, M.\ (2023). \\
\emph{High-Dimensional Private Empirical Risk Minimization by Greedy Coordinate Descent.} AISTATS 2023.
\end{frame}

\end{document}